% Stage 8 restored fragment — 2.4.2 Outgoing-wave selection (archive lines 5231-5455)
\subsection{Outgoing-wave selection}
\label{sec:ch242}
Finite computational domains terminate open-region problems artificially. On a physical compact
range, anechoic absorbers play the same role as absorbing boundary conditions in equivalent analytical models:
they mimic terminations that let export escape instead of reflecting into the near field
\cite{hansen1988}. Silver--M\"uller and related traces approximate outgoing selection
on meshes used for horns, cavities, and reflector feeds
\cite{stratton1941}. A finite-element model employing inadequate
absorbing boundaries permits artificial reflections that contaminate field amplitudes and
reconstructed radiation characteristics---the same standing-wave signatures that appear when chamber
absorption is imperfect on a physical range \cite{hansen1988}. Absorbing-boundary
parameters are therefore a radiation-physics choice about exported energy, not a mesh-only tuning
parameter \cite{stratton1941}. The operator $\mathcal{O}_{\mathrm{out}}$ records that choice on
$\mathcal{S}_\omega$ before $\Pi_{\mathrm{rad}}$ and $\mathcal{R}_\omega$ attach pattern meters.

The outgoing-wave selector $\mathcal{O}_{\mathrm{out}}$ implements Sommerfeld completion
Eq.~\eqref{eq:ch2.sommerfeld-outgoing-class} as an operator on $\mathcal{F}_\Omega$, on
harmonic Green kernels, and on the resolvent branch of $\mathcal{L}_\omega^{-1}$
\cite{stratton1941}. The composition relation
Eq.~\eqref{eq:ch2.rad-projection-composition} from Subsection~\ref{sec:ch224} named
$\mathcal{O}_{\mathrm{out}}$ without constructing it; the definitions below supply that
operator and its links to $\Pi_{\mathrm{prop}}$, $\Pi_{\mathrm{rad}}$, and retarded
Green functions Eq.~\eqref{eq:ch2.retarded-harmonic-kernel-theorem}. Sommerfeld field
conditions
Eqs.~\eqref{eq:ch2.sommerfeld-field-condition}--\eqref{eq:ch2.sommerfeld-outgoing-class}
are presupposed; uniqueness proofs belong to Subsection~\ref{sec:ch243}.
Assumptions. Exterior regions use vacuum far-zone constants with \(k_0>0\) real in the locked
harmonic band, compact \(\operatorname{supp}\Jimp\), and time-independent
\(\mathcal{B}_\Gamma\) including outgoing radiation condition at infinity or on
Eq.~\eqref{eq:ch2.silver-muller-trace-approx} at large finite \(R\)
\cite{stratton1941}. Propagating and evanescent projectors
Eqs.~\eqref{eq:ch2.spectral-projectors}--\eqref{eq:ch2.spectral-projector-algebra}
from Subsection~\ref{sec:ch224} are fixed on \(\widehat{\mathcal{F}}_\Omega\)
\cite{harrington2001}. Incoming and advanced branches satisfying
Eq.~\eqref{eq:ch2.advanced-kernel-violation} are excluded.
Let \(\mathcal{F}_{\mathrm{in}}\) denote the incoming counterpart of
\(\mathcal{F}_{\mathrm{SF}}\) obtained by reversing outgoing phase in the far zone.
The outgoing-wave selector is the idempotent map
\begin{equation}
\label{eq:ch2.outgoing-selector-def}
\mathcal{O}_{\mathrm{out}}:\mathcal{F}_\Omega\to\mathcal{F}_{\mathrm{SF}},
\qquad
\mathcal{O}_{\mathrm{out}}[\mathbf{u}_\omega]
=
\mathbf{u}_{\mathrm{out}},
\end{equation}
where \(\mathbf{u}_{\mathrm{out}}\) is the unique member of
\(\mathcal{F}_{\mathrm{SF}}\) whose propagating spectral content matches the outgoing
\(k_z\) branch of Eq.~\eqref{eq:ch2.longitudinal-wavenumber-split} and whose
incoming content is annihilated \cite{stratton1941}. Equation
\eqref{eq:ch2.outgoing-selector-def} is the locked definition for operator
realization of Sommerfeld completion: it acts \emph{before} far-zone measurement and
\emph{after} volume completion by \(\mathcal{L}_\omega\), selecting the physical
resolvent branch assumed in Eq.~\eqref{eq:ch2.resolvent-green-equivalence}
\cite{stratton1941}. On \(\mathcal{F}_{\mathrm{SF}}\),
\begin{equation}
\label{eq:ch2.outgoing-idempotency}
\mathcal{O}_{\mathrm{out}}^2=\mathcal{O}_{\mathrm{out}},
\qquad
\mathcal{O}_{\mathrm{out}}\big|_{\mathcal{F}_{\mathrm{SF}}}=\mathrm{id},
\end{equation}
while incoming states satisfy \(\mathcal{O}_{\mathrm{out}}[\mathbf{u}_{\mathrm{in}}]=\mathbf{0}\)
when extended by continuity from the far zone \cite{harrington2001}. Physically,
\(\mathcal{O}_{\mathrm{out}}\) removes incoming spherical and plane-wave components
that satisfy Maxwell equations but violate retarded causality
\cite{jackson1999}.
On homogeneous regions, outgoing selection in lateral wavenumber \((k_x,k_y)\) uses
the longitudinal root from Eq.~\eqref{eq:ch2.longitudinal-wavenumber-split}:
\begin{equation}
\label{eq:ch2.outgoing-plane-wave-criterion}
k_z=
\begin{cases}
+\sqrt{k_0^2-k_t^2}, & k_t\le k_0\ \text{(outgoing propagating)},\\[4pt]
-\jj\,\alpha(k_t), & k_t>k_0\ \text{(decaying evanescent)},
\end{cases}
\end{equation}
with the sign of the propagating root fixed by exterior completion and the locked
\(\exp(-\jj\omega t)\) convention \cite{stratton1941}. Equation
\eqref{eq:ch2.outgoing-plane-wave-criterion} is the spectral action of
\(\mathcal{O}_{\mathrm{out}}\) on each plane-wave mode; reversing
\(k_z\to -k_z\) on the propagating sheet produces incoming content annihilated by
\(\mathcal{O}_{\mathrm{out}}\) \cite{stratton1941}. Evanescent modes are unchanged
in decay direction but remain outside \(\Pi_{\mathrm{rad}}\) through
Eq.~\eqref{eq:ch2.rad-null-on-evanescent} \cite{harrington2001}. The propagating
projector \(\Pi_{\mathrm{prop}}\) from Subsection~\ref{sec:ch224} retains all
outgoing and incoming propagating modes; \(\mathcal{O}_{\mathrm{out}}\) acts only
after \(\Pi_{\mathrm{prop}}\) in the radiation chain
Eq.~\eqref{eq:ch2.rad-projection-composition} \cite{stratton1941}.
\begin{figure}[t]
\centering
\ChOneFigBlock{%
\begin{tikzpicture}[ch1fig, node distance=5mm and 7mm]
 \node[ch1box] (f) {$\mathcal{F}_\Omega$};
 \node[ch1box, right=7mm of f] (pp) {$\Pi_{\mathrm{prop}}$};
 \node[ch1box, right=7mm of pp] (oo) {$\mathcal{O}_{\mathrm{out}}$};
 \node[ch1box, right=7mm of oo] (pr) {$\Pi_{\mathrm{rad}}$};
 \node[ch1box, below=5mm of f] (ev) {$\Pi_{\mathrm{ev}}$};
 \draw[ch1flow] (f) -- (pp) -- (oo) -- (pr);
 \draw[ch1flow] (f.south) -- (f.south |- ev.north);
 \draw[ch1flow] (f.south |- ev.north) -- (ev.north);
\end{tikzpicture}%
}
\caption{Outgoing selection Eq.~\eqref{eq:ch2.outgoing-selector-def} sits between
propagating projection and radiation export:
\(\Pi_{\mathrm{rad}}=\mathcal{O}_{\mathrm{out}}\circ\Pi_{\mathrm{prop}}\).}
\label{fig:ch2.outgoing-selector-measurement chain}
\end{figure}
~\eqref{eq:ch2.rad-projection-composition} with explicit operator order. Regime
qualification Eq.~\eqref{eq:ch2.regime-operator-screen} applies after
\(\Pi_{\mathrm{rad}}\), not at the \(\mathcal{O}_{\mathrm{out}}\) stage
Green-function branch selection. Harmonic kernels from Section~\ref{sec:ch23}---with observation coordinate $\mathbf{r}$ and
source coordinate $\mathbf{r}'$ in $G$ and $\overline{\mathbf{G}}_{E}$---admit incoming and
outgoing analytic continuations. The outgoing dyadic branch is
\begin{equation}
\label{eq:ch2.outgoing-green-branch}
\overline{\mathbf{G}}_{E,\omega}^{\mathrm{(out)}}
=
\mathcal{O}_{\mathrm{out}}\,\overline{\mathbf{G}}_{E,\omega},
\qquad
\overline{\mathbf{G}}_{E,\omega}^{\mathrm{(out)}}
\ \text{satisfies}\ 
\text{Eq.~\eqref{eq:ch2.sommerfeld-dyadic-asymptotic}},
\end{equation}
with \(\mathcal{O}_{\mathrm{out}}\) acting on the observation-point dependence
\cite{stratton1941}. Equation~\eqref{eq:ch2.outgoing-green-branch}
identifies the kernel used in Eq.~\eqref{eq:ch2.electric-dyadic-convolution} on
exterior problems with the retarded symbol of
Eq.~\eqref{eq:ch2.retarded-harmonic-kernel-theorem} \cite{stratton1941}. Incoming
kernels \(\overline{\mathbf{G}}_{E,\omega}^{\mathrm{(in)}}\) are related by
\(k_0\to -k_0\) on the propagating Riemann sheet and produce fields outside
\(\mathcal{F}_{\mathrm{SF}}\) \cite{jackson1999}. Subsection~\ref{sec:ch242} does
not re-tabulate Eq.~\eqref{eq:ch2.free-space-dyadic-green}; it fixes the sheet on
which that formula is evaluated \cite{harrington2001}. The physical resolvent branch is
\begin{equation}
\label{eq:ch2.outgoing-resolvent-selection}
\mathcal{L}_\omega^{-1}\big|_{\mathrm{phys}}
=
\mathcal{O}_{\mathrm{out}}\circ\mathcal{L}_\omega^{-1},
\end{equation}
understood on \(\mathcal{D}(\mathcal{L}_\omega)\) with outgoing \(\mathcal{B}_\Gamma\)
\cite{stratton1941}. The solution map on impressed currents becomes
\begin{equation}
\label{eq:ch2.outgoing-solution-operator}
\mathcal{S}_\omega[\Jimp]
=
\mathcal{O}_{\mathrm{out}}\,
\Big(-\jj\omega\mu_0\int_{\Omega}
\overline{\mathbf{G}}_{E,\omega}^{\mathrm{(out)}}(\mathbf{r},\mathbf{r}')\cdot\Jimp(\mathbf{r}')\,d^{3}r'
\;+\;\E_{\mathrm{hom}}\Big),
\end{equation}
with regularization from Subsection~\ref{sec:ch232} when \(\mathbf{r}\) approaches
\(\operatorname{supp}\Jimp\) \cite{harrington2001}. Equations
\eqref{eq:ch2.outgoing-resolvent-selection} and~\eqref{eq:ch2.outgoing-solution-operator}
realize Eq.~\eqref{eq:ch2.exterior-uniqueness-hypothesis} once
Subsection~\ref{sec:ch243} proves uniqueness on the selected branch
\cite{stratton1941}. Causality from Subsection~\ref{sec:ch233} is consistent with
Eq.~\eqref{eq:ch2.outgoing-green-branch} but not sufficient when harmonic solvers
bypass time-domain construction \cite{harrington2001}. The volume integral in
Eq.~\eqref{eq:ch2.outgoing-solution-operator} sums impressed current contributions at each
source point \(\mathbf{r}'\in\Omega\) against the outgoing dyadic evaluated at observation
point \(\mathbf{r}\); \(\E_{\mathrm{hom}}\) carries any exterior illumination not generated
by \(\Jimp\) and must vanish for pure radiation problems \cite{stratton1941}.
On propagating content, outgoing selection precedes radiation export in the measurement sequence
fixed in Section~\ref{sec:ch24}: $\mathcal{R}_\omega=\Pi_{\mathrm{rad}}\circ\mathcal{S}_\omega$ with
$\Pi_{\mathrm{rad}}=\mathcal{O}_{\mathrm{out}}\circ\Pi_{\mathrm{prop}}$ \cite{stratton1941}. On
propagating content,
\begin{equation}
\label{eq:ch2.outgoing-radiation-refinement}
\Pi_{\mathrm{rad}}
=
\mathcal{O}_{\mathrm{out}}\circ\Pi_{\mathrm{prop}},
\qquad
\mathcal{R}_\omega
=
\Pi_{\mathrm{rad}}\circ\mathcal{S}_\omega,
\end{equation}
restating Eqs.~\eqref{eq:ch2.rad-projection-composition} and~\eqref{eq:ch2.radiation-operator-def}
with the locked definition of \(\mathcal{O}_{\mathrm{out}}\)
\cite{stratton1941}. Radiative subspace
\(\mathcal{F}_{\mathrm{rad}}\subset\mathcal{F}_{\mathrm{SF}}\) consists of fields
with nonzero far-zone transverse amplitude Eq.~\eqref{eq:ch2.far-zone-transverse-radiation};
\(\mathcal{O}_{\mathrm{out}}\) maps into \(\mathcal{F}_{\mathrm{SF}}\) but not
automatically into \(\mathcal{F}_{\mathrm{rad}}\) when evanescent content dominates
near sources \cite{stratton1941}. Transport reporting therefore requires
Eq.~\eqref{eq:ch2.regime-operator-screen} after
Eq.~\eqref{eq:ch2.outgoing-radiation-refinement} \cite{harrington2001}.
Outgoing selection on a compact range is implemented twice in practice: absorbers suppress
multipath reflections that would mimic incoming waves, and harmonic solvers fix the propagating
$k_z$ branch of Eq.~\eqref{eq:ch2.outgoing-plane-wave-criterion} on exterior meshes
\cite{hansen1988}. When either layer fails, $\mathcal{O}_{\mathrm{out}}$ acts on
the wrong field orbit even though volume Maxwell rows still vanish \cite{harrington2001}.
Incoming--outgoing split. Define \(\mathcal{O}_{\mathrm{in}}\) analogously on \(\mathcal{F}_{\mathrm{in}}\).
On the propagating subspace,
\begin{equation}
\label{eq:ch2.incoming-outgoing-complement}
\mathrm{id}
=
\mathcal{O}_{\mathrm{out}}+\mathcal{O}_{\mathrm{in}}
\quad\text{on}\ \Pi_{\mathrm{prop}}\mathcal{F}_\Omega,
\qquad
\mathcal{O}_{\mathrm{out}}\mathcal{O}_{\mathrm{in}}=0,
\end{equation}
with evanescent and non-propagating components orthogonal to both when decay
direction is fixed by Eq.~\eqref{eq:ch2.outgoing-plane-wave-criterion}
\cite{stratton1941}. Equation~\eqref{eq:ch2.incoming-outgoing-complement}
explains branch degeneracy in unconstrained Helmholtz solvers: both
\(\mathcal{O}_{\mathrm{out}}[\mathbf{u}]\) and
\(\mathcal{O}_{\mathrm{in}}[\mathbf{u}]\) satisfy volume equations unless one is
selected by boundary or radiation completion \cite{jackson1999}.
Equations~\eqref{eq:ch2.outgoing-selector-def}--\eqref{eq:ch2.outgoing-solution-operator}
require linear media, fixed \(\omega\), and outgoing \(\mathcal{B}_\Gamma\) on
exterior problems. They fail on closed PEC cavities where
Eq.~\eqref{eq:ch1.outer-boundary-classes} replaces Sommerfeld completion, and on
lossless open resonators where discrete spectra require quotienting
Eq.~\eqref{eq:ch2.coercivity-quotient} \cite{stratton1941}.
Absorbing layers approximate \(\mathcal{O}_{\mathrm{out}}\) but introduce spurious
modes unless tuned to Eq.~\eqref{eq:ch2.sommerfeld-field-condition} at large distance
\cite{harrington2001}.
